\documentclass[12pt,a4paper]{article}

% =====================================================================
% PACKAGES
% =====================================================================
\usepackage[utf8]{inputenc}
\usepackage[T1]{fontenc}
\usepackage[german]{babel}
\usepackage{amsmath, amssymb, amsfonts}
\usepackage{geometry}
\usepackage{graphicx}
\usepackage{hyperref}
\usepackage{xcolor}
\usepackage{titlesec}
\usepackage{fancyhdr}
\usepackage{booktabs}
\usepackage{enumitem}
\usepackage{authblk}
\usepackage{natbib}
\usepackage{setspace}
\usepackage{tcolorbox}

% =====================================================================
% LAYOUT
% =====================================================================
\geometry{left=2.5cm, right=2.5cm, top=3cm, bottom=3cm}
\onehalfspacing

% =====================================================================
% COLORS & LINKS
% =====================================================================
\definecolor{darkblue}{RGB}{0,51,102}
\definecolor{warnorange}{RGB}{204,102,0}
\hypersetup{
    colorlinks=true,
    linkcolor=darkblue,
    citecolor=darkblue,
    urlcolor=darkblue
}

% =====================================================================
% SECTION FORMATTING
% =====================================================================
\titleformat{\section}{\Large\bfseries\color{darkblue}}{\thesection}{1em}{}
\titleformat{\subsection}{\large\bfseries\color{darkblue!80}}{\thesubsection}{1em}{}

% =====================================================================
% HEADER & FOOTER
% =====================================================================
\pagestyle{fancy}
\fancyhf{}
\fancyhead[R]{\small\color{gray}SOTA: Actionable Theories \& Causal Constructs}
\fancyfoot[C]{\thepage}
\renewcommand{\headrulewidth}{0.4pt}

% =====================================================================
% DOCUMENT INFO
% =====================================================================
\title{
{\LARGE\bfseries State of the Art Report (SOTA)}\\[0.4cm]
{\Large Building Actionable Theories -- The Role of Causal Constructs (Dietz, 2026)}\\[0.2cm]
{\large und die Implementation in BEATRIX / BCM 2.2}
}

\author[1]{Andrea Fehr}
\author[1]{Gerhard Fehr}
\affil[1]{FehrAdvice \& Partners AG, Zürich}
\date{Dezember 2025}

% =====================================================================
% DOCUMENT
% =====================================================================
\begin{document}

\maketitle

\begin{abstract}
Dieser Bericht untersucht die theoretische und methodologische Bedeutung von Dietz (2026) \textit{Building Actionable Theories: The Role of Causal Constructs} im Kontext der Entwicklung kausaler, manipulationsfähiger Modelle in den Sozial- und Verhaltenswissenschaften.
Das Paper schlägt vor, dass Theorien nur dann handlungsleitend sind, wenn ihre Konstrukte manipulierbar, messbar und in überprüfbare Kausalnetzwerke eingebettet sind.
Diese Arbeit analysiert, inwiefern das Behavioral Change Model (BCM 2.2) und dessen Implementierung in BEATRIX diesen konzeptionellen Ansatz operationalisieren -- einschließlich einer kritischen Diskussion der \textbf{Grenzen} dieser Operationalisierung, insbesondere hinsichtlich des Lab-Field Gaps und der Unterscheidung zwischen Modell-Manipulation und Welt-Manipulation.
\end{abstract}

% =====================================================================
\section{Einleitung: Theorie und Handlungsfähigkeit}
% =====================================================================

In der sozialwissenschaftlichen Theoriediskussion besteht seit Jahrzehnten die Spannung zwischen erklärenden und handlungsleitenden Ansätzen.
Während klassische Modelle darauf abzielen, Phänomene zu beschreiben oder zu erklären, bleibt häufig offen, wie dieses Wissen in konkretes Handeln überführt werden kann.
Dietz (2026) adressiert dieses Problem, indem er die Bedingung formuliert, unter der Theorien als \textit{actionable} gelten -- also als fähig, Interventionen zu leiten oder Verhalten gezielt zu verändern.
Er positioniert seine Arbeit im Schnittpunkt zwischen Kausalitätsforschung (Woodward, 2003; Pearl, 2009) und Theoriekonstruktion (Weick, 1989; Sutton \& Staw, 1995) und schlägt ein Konzept vor, das Kausalität als Grundlage von Handlungsfähigkeit versteht.

% =====================================================================
\section{Stand der Forschung}
% =====================================================================

\subsection{Die erklärungsorientierte Phase (1960--1990)}

Frühe Theoriebildung in Ökonomie, Psychologie und Management zielte auf Erklärung und Generalisierung ab.
Arbeiten wie Simon (1962) zu begrenzter Rationalität oder Bandura (1977) zur Selbstwirksamkeit konzentrierten sich auf beobachtbare Muster, nicht auf Eingriffsmöglichkeiten.
Das Wissen blieb beschreibend; Theorien erklärten Phänomene, ohne sie gezielt reproduzieren oder verändern zu können.

\subsection{Die kontextuelle Wende (1990--2015)}

In der Folge rückten System- und Kontextfaktoren in den Vordergrund.
Eisenhardt (1999) und Giddens (2000) betonten die Bedeutung situativer und struktureller Bedingungen für menschliches Verhalten.
Dennoch blieb die Frage der gezielten Manipulierbarkeit theoretischer Variablen weitgehend unbeantwortet.
Die Theorien blieben diagnostisch, nicht interventionistisch.

\subsection{Von Erklärung zu Intervention (2015--2025)}

Seit der Etablierung evidenzbasierter Ansätze (Rousseau, 2018; Edmondson, 2021) wuchs das Interesse an \textit{actionable theories}.
Dietz (2026) greift diesen Trend auf und entwickelt eine präzise Definition, wann theoretische Konstrukte manipulierbar und empirisch steuerbar sind.
Damit wird die Lücke zwischen Theorie und Anwendung methodologisch geschlossen.

% =====================================================================
\section{Der Beitrag von Dietz (2026): Causal Constructs}
% =====================================================================

Dietz unterscheidet zwischen drei Klassen theoretischer Kausalität:

\begin{center}
\begin{tabular}{llll}
\toprule
\textbf{Typ} & \textbf{Beschreibung} & \textbf{Beispiel} & \textbf{Actionability}\\
\midrule
Observe(X) & Beobachtung von Zusammenhängen & Fairness korreliert mit Zufriedenheit & Nein \\
Select(X) & Auswahl- oder Selbstselektion & Selbstwirksame Personen wählen Aufgaben & Eingeschränkt \\
Manipulate(X) & Aktive Veränderung erzeugt Wirkung & Erhöhte Feedbackfrequenz steigert Motivation & Ja \\
\bottomrule
\end{tabular}
\end{center}

Nur der dritte Typ -- \textit{Manipulate(X)} -- erfüllt die Bedingung einer handlungsleitenden Theorie.
Dietz definiert Theorien als \textit{actionable}, wenn sie mindestens ein Konstrukt enthalten, das gezielt verändert werden kann und dessen Veränderung eine überprüfbare Wirkung entfaltet.

% =====================================================================
\section{Causal Construct Criteria}
% =====================================================================

Dietz formuliert vier Bedingungen, die erfüllt sein müssen, damit ein theoretisches Konstrukt als kausal handlungsrelevant gilt:

\begin{itemize}[leftmargin=1.5em]
    \item \textbf{Distinctiveness:} Das Konstrukt ist klar definiert und abgegrenzt.
    \item \textbf{Nomological Embeddedness:} Es ist in ein überprüfbares Kausalnetzwerk eingebettet.
    \item \textbf{Manipulability:} Das Konstrukt kann konzeptionell oder experimentell verändert werden.
    \item \textbf{Operationalizability:} Das Konstrukt ist messbar und quantifizierbar.
\end{itemize}

Diese vier Kriterien schaffen die Grundlage für die Übertragung theoretischer Modelle in rechnerische oder experimentelle Systeme.

% =====================================================================
\section{Theoretischer Rahmen}
% =====================================================================

Dietz' Ansatz basiert auf der interventionistischen Kausalitätslogik.
Woodward (2003) versteht Kausalität als Möglichkeit, durch Veränderung einer Variable eine andere zu beeinflussen.
Pearl (2009) formuliert diese Logik mathematisch über den sogenannten \textit{do-Operator}, der Interventionen modelliert.
Bunge (2015) erweitert diese Sicht, indem er die Mechanismen betont, durch die Ursachen Wirkungen erzeugen.

Dietz verbindet diese Perspektiven mit Theoriearbeit und argumentiert, dass nur Theorien, die explizit manipulierbare Variablen enthalten, empirisch fruchtbar sind.

% =====================================================================
\section{Methodische Konsequenzen}
% =====================================================================

Dietz schlägt vor, Theorien so zu strukturieren, dass sie in Simulations- oder Modellarchitekturen überführbar sind.
Dazu müssen Variablen parametrisiert und ihre Wechselwirkungen formal beschreibbar sein.
Diese Forderung führt zu einer formalen Theorielogik, in der erklärende und manipulative Komponenten systematisch getrennt sind.
Somit wird Theoriebildung selbst zu einem Designprozess, der Eingriffsmöglichkeiten explizit abbildet.

% =====================================================================
\section{Integration in BEATRIX und BCM 2.2}
% =====================================================================

\subsection{Zwei-Ebenen-Architektur}

Das BEATRIX-System nutzt eine zweistufige Struktur:

\begin{center}
\begin{tabular}{lll}
\toprule
\textbf{Schicht} & \textbf{Funktion} & \textbf{Charakteristik}\\
\midrule
LLM-Schicht & Semantisches Parsing von Sprache & probabilistisch\\
BCM-Schicht & Berechnung der Modellstruktur & strukturiert (mit Unsicherheit)\\
\bottomrule
\end{tabular}
\end{center}

\begin{tcolorbox}[colback=yellow!5!white, colframe=warnorange, title=Präzisierung: ``Deterministisch'' vs. ``Strukturiert'']
Eine frühere Version dieses Dokuments bezeichnete die BCM-Schicht als ``deterministisch''. Dies ist eine Vereinfachung. BCM propagiert Unsicherheit durch Monte-Carlo-Simulation und liefert Konfidenzintervalle, nicht Punktschätzungen. Die korrekte Charakterisierung ist \textbf{strukturiert mit quantifizierter Unsicherheit}.
\end{tcolorbox}

Diese Trennung entspricht der inhaltlichen Logik von Dietz.
Die sprachliche Interpretation (LLM) liefert Parameter, während das mathematische Modell (BCM) die kausale Struktur berechnet.

\subsection{Bezug zur Theorie von Dietz}

\begin{center}
\begin{tabular}{ll}
\toprule
\textbf{Dietz (2026)} & \textbf{BEATRIX / BCM}\\
\midrule
Manipulable Construct & Exogene Parameter (z.B. $\alpha, \beta, \rho, \eta$)\\
Nomological Network & Modellgleichungen im BCM\\
Operationalizability & Rechenbare Nutzenfunktionen\\
Producing Cause & Veränderung von Verhalten und Nutzen\\
\bottomrule
\end{tabular}
\end{center}

Diese Zuordnung verdeutlicht, dass BEATRIX die von Dietz beschriebenen logischen Bedingungen technisch abbildet.
\textbf{Jedoch ist eine kritische Differenzierung erforderlich}, die im folgenden Abschnitt ausgeführt wird.

% =====================================================================
\section{Kritische Differenzierung: Modell-Manipulation vs. Welt-Manipulation}
% =====================================================================

\subsection{Zwei Bedeutungen von ``Manipulierbarkeit''}

Dietz' Kriterium der \textit{Manipulability} kann auf zwei Ebenen interpretiert werden:

\begin{center}
\begin{tabular}{p{5.5cm}p{6.5cm}}
\toprule
\textbf{Modell-Manipulation} & \textbf{Welt-Manipulation}\\
\midrule
Parameter können im Modell geändert werden & Interventionen ändern Parameter in der Realität\\
$\lambda$ kann von 2.0 auf 2.5 gesetzt werden & Ein Nudge \textit{verursacht} höhere Loss Aversion\\
\textit{What-if}-Analyse & \textit{Kausaler Effekt}\\
Simulation & Intervention\\
\textbf{BCM leistet dies: }\checkmark & \textbf{BCM ermöglicht dies, garantiert es aber nicht}\\
\bottomrule
\end{tabular}
\end{center}

\begin{tcolorbox}[colback=blue!5!white, colframe=darkblue, title=Zentrale Unterscheidung]
BCM/BEATRIX ermöglicht \textbf{Modell-Manipulation}: Die Konsequenzen von Parameteränderungen können berechnet werden.

Ob eine reale Intervention tatsächlich die gewünschte Parameteränderung bewirkt, ist eine \textbf{empirische Frage}, die das Modell nicht beantworten kann.
\end{tcolorbox}

\subsection{Implikation für Dietz' Kriterien}

\begin{center}
\begin{tabular}{llc}
\toprule
\textbf{Kriterium} & \textbf{BCM/BEATRIX Status} & \textbf{Erfüllung}\\
\midrule
Distinctiveness & 115 explizite Axiome, klare Parameterdefinitionen & \checkmark\\
Nomological Embeddedness & Formale Modellgleichungen, 5-Ebenen-Hierarchie & \checkmark\\
Manipulability (Modell) & Parameter als exogene Inputs änderbar & \checkmark\\
Manipulability (Welt) & Abhängig von Interventionswirksamkeit & $\sim$\\
Operationalizability & Quantifizierte Parameter mit Konfidenzintervallen & \checkmark\\
\bottomrule
\end{tabular}
\end{center}

BCM erfüllt Dietz' Kriterien \textbf{auf Modellebene} vollständig. Die Übertragung auf \textbf{Welt-Manipulation} erfordert zusätzliche empirische Validierung.

% =====================================================================
\section{Der Lab-Field Gap: Empirische Grenzen der Manipulierbarkeit}
% =====================================================================

\subsection{Das Problem}

Ein zentrales Ergebnis der Verhaltensökonomie der letzten Jahre ist der sogenannte \textit{Lab-Field Gap}: Effekte, die in kontrollierten Laborexperimenten gefunden werden, replizieren im Feld oft mit deutlich geringerer Stärke.

\textbf{Zentrale Befunde:}
\begin{itemize}
    \item DellaVigna \& Linos (2022): Nudge-Effekte im Feld betragen \textbf{10--20\% der Laboreffekte}
    \item Allcott (2015): Soziale Vergleichs-Nudges zeigen starke Heterogenität
    \item Hummel \& Maedche (2019): Meta-Analyse zeigt systematische Überschätzung von Nudge-Wirksamkeit
\end{itemize}

\subsection{BCM-Axiom MA-EPI-38: Evidenzhierarchie}

Das BCM adressiert dieses Problem explizit durch das Axiom MA-EPI-38, das eine Evidenzhierarchie für Parameter definiert:

\begin{center}
\begin{tabular}{clp{7cm}}
\toprule
\textbf{Stufe} & \textbf{Evidenzqualität} & \textbf{Beschreibung}\\
\midrule
$\star\star\star\star$ & Höchste & Konvergenz von Labor-RCT + Feld-RCT + Natürlichem Experiment\\
$\star\star\star$ & Hohe & Replizierte Labor-RCTs mit theoretischer Fundierung\\
$\star\star$ & Mittlere & Einzelne Labor-RCTs oder nicht-replizierte Feldstudien\\
$\star$ & Niedrige & Beobachtungsstudien oder theoretische Ableitung\\
\bottomrule
\end{tabular}
\end{center}

\begin{tcolorbox}[colback=red!5!white, colframe=red!75!black, title=Implikation für Actionability]
Die meisten BCM-Parameter befinden sich auf Evidenzstufe $\star\star$ oder $\star\star\star$.

\textbf{Das bedeutet}: Die \textit{Manipulierbarkeit} im Sinne von Dietz ist für viele Parameter empirisch nur unter Laborbedingungen nachgewiesen. Die Übertragung auf Feldkontexte erfordert Vorsicht und kontextspezifische Validierung.
\end{tcolorbox}

\subsection{Konsequenz für BEATRIX-Prognosen}

BEATRIX-Prognosen sollten als \textbf{kalibrierte Schätzungen unter Unsicherheit} verstanden werden, nicht als präzise Vorhersagen:

\begin{center}
\begin{tabular}{p{5cm}p{7cm}}
\toprule
\textbf{BEATRIX sagt} & \textbf{Korrekte Interpretation}\\
\midrule
``4.2\% Churn [2.8\%, 5.9\%]'' & ``Basierend auf verfügbarer Evidenz und BCM-Struktur ist 4.2\% der wahrscheinlichste Wert, mit erheblicher Unsicherheit''\\
\midrule
``Nudge X reduziert Churn um 1.5pp'' & ``Unter Laborbedingungen reduziert Nudge X Churn um 1.5pp. Feldeffekt vermutlich geringer''\\
\bottomrule
\end{tabular}
\end{center}

% =====================================================================
\section{Was BCM/BEATRIX leistet -- und was nicht}
% =====================================================================

\subsection{Stärken (wo BCM Dietz' Vision erfüllt)}

\begin{enumerate}
    \item \textbf{Strukturierte Theoriebildung}: 115 explizite Axiome machen Annahmen transparent und kritisierbar
    \item \textbf{Nomologisches Netzwerk}: Formale Gleichungen verbinden Parameter mit Outcomes
    \item \textbf{Operationalisierung}: Alle Konstrukte sind quantifiziert und berechenbar
    \item \textbf{Interventions-Taxonomie}: 8 Interventionstypen mit expliziten Mechanismen
    \item \textbf{What-if-Analyse}: Konsequenzen von Parameteränderungen sind simulierbar
    \item \textbf{Dokumentation}: Vollständige Reasoning-Chain für Nachvollziehbarkeit
    \item \textbf{Lernfähigkeit}: Outcome-Feedback kann zur Kalibrierung genutzt werden
\end{enumerate}

\subsection{Grenzen (wo Vorsicht geboten ist)}

\begin{enumerate}
    \item \textbf{Modell $\neq$ Realität}: Parameteränderungen im Modell sind nicht automatisch in der Welt erreichbar
    \item \textbf{Lab-Field Gap}: Interventionseffekte sind im Feld typischerweise 10--20\% der Laboreffekte
    \item \textbf{Kein formaler do-Kalkül}: BCM implementiert nicht Pearls do-Operator; die kausale Struktur ist konzeptuell, nicht identifiziert
    \item \textbf{Evidenzheterogenität}: Parameter haben unterschiedliche Evidenzqualität ($\star$ bis $\star\star\star\star$)
    \item \textbf{Kontextabhängigkeit}: Effekte variieren stark nach Domäne, Population und Implementierung
    \item \textbf{LLM-Schätzunsicherheit}: Die LLM-basierte Parameterschätzung fügt eine zusätzliche Unsicherheitsebene hinzu
\end{enumerate}

\subsection{Die korrekte Positionierung}

\begin{tcolorbox}[colback=green!5!white, colframe=green!75!black, title=BCM/BEATRIX als ``Actionable Theory''-Framework]
BCM/BEATRIX ist kein Orakel, das präzise Verhaltensvorhersagen liefert.

Es ist ein \textbf{Framework für strukturiertes Denken über Verhaltensänderung}, das:
\begin{itemize}
    \item Annahmen \textbf{explizit} macht (statt implizit)
    \item Unsicherheit \textbf{quantifiziert} (statt ignoriert)
    \item Mechanismen \textbf{spezifiziert} (statt vage beschreibt)
    \item Interventionen \textbf{systematisiert} (statt ad-hoc wählt)
    \item Lernen \textbf{ermöglicht} (statt Fehler zu wiederholen)
\end{itemize}

Im Sinne von Dietz (2026) ist BCM \textit{actionable} -- aber die \textit{action} erfordert empirische Validierung im spezifischen Kontext.
\end{tcolorbox}

% =====================================================================
\section{Diskussion}
% =====================================================================

Dietz' Ansatz liefert eine präzise Definition, wann Theorien handlungsfähig sind.
Er ermöglicht eine systematische Verbindung zwischen theoretischem Wissen und rechnerischer Umsetzung.
BCM/BEATRIX illustriert die Umsetzbarkeit dieses Denkens, indem es kausale Strukturen explizit formalisiert.

\textbf{Drei Differenzierungen sind jedoch wesentlich:}

\begin{enumerate}
    \item \textbf{Modell-Manipulation $\neq$ Welt-Manipulation}: BCM ermöglicht strukturierte What-if-Analysen. Ob Interventionen die gewünschten Parameteränderungen bewirken, ist empirisch zu prüfen.

    \item \textbf{Lab-Field Gap}: Die verhaltensökonomische Forschung zeigt, dass Interventionseffekte im Feld systematisch geringer sind als im Labor. BCM adressiert dies durch die Evidenzhierarchie (MA-EPI-38), aber Anwender müssen diese Unsicherheit berücksichtigen.

    \item \textbf{Strukturierte Unsicherheit statt Punktvorhersagen}: BCM liefert keine deterministischen Prognosen, sondern Schätzungen mit Konfidenzintervallen. Diese Unsicherheit ist ein Feature, kein Bug -- sie reflektiert die tatsächliche Erkenntnislage.
\end{enumerate}

Offen bleibt, wie weit diese Prinzipien empirisch verallgemeinerbar sind und welche Grenzen die formale Manipulierbarkeit in komplexen sozialen Systemen hat.

% =====================================================================
\section{Fazit}
% =====================================================================

Der Beitrag von Dietz (2026) stellt einen methodologischen Bezugspunkt für die Diskussion über \textit{actionable theories} dar.
Er formuliert klare Bedingungen, unter denen Theorien manipulierbar und überprüfbar werden.
BCM 2.2 und BEATRIX lassen sich als technische Umsetzungen dieser Bedingungen interpretieren -- \textbf{auf Modellebene}.

Die Analyse zeigt:
\begin{itemize}
    \item BCM erfüllt Dietz' Kriterien (Distinctiveness, Nomological Embeddedness, Operationalizability) vollständig
    \item Manipulability ist differenziert zu betrachten: Modell-Manipulation ist gegeben, Welt-Manipulation erfordert empirische Validierung
    \item Der Lab-Field Gap ist eine systematische Herausforderung, die BCM durch die Evidenzhierarchie adressiert, aber nicht eliminiert
    \item BEATRIX ist ein Framework für strukturiertes Denken, kein Prognose-Orakel
\end{itemize}

\begin{tcolorbox}[colback=blue!5!white, colframe=darkblue, title=Schlussfolgerung]
BCM/BEATRIX macht Verhaltenstheorie \textit{actionable} im Sinne von Dietz -- nicht durch Garantie präziser Vorhersagen, sondern durch \textbf{Strukturierung}, \textbf{Quantifizierung} und \textbf{Transparenz} von Annahmen, Mechanismen und Unsicherheiten.

Die praktische Handlungsfähigkeit entsteht nicht aus dem Modell allein, sondern aus der \textbf{Kombination} von strukturierter Theorie, empirischer Validierung und kontextspezifischer Anpassung.
\end{tcolorbox}

% =====================================================================
\section*{Literatur}
% =====================================================================

\begin{thebibliography}{99}

\bibitem[Allcott, 2015]{allcott2015}
Allcott, H. (2015). Site Selection Bias in Program Evaluation. \textit{Quarterly Journal of Economics, 130}(3), 1117--1165.

\bibitem[Argyris \& Schön, 1996]{argyris1996}
Argyris, C., \& Schön, D. (1996). \textit{Theory in Practice: Increasing Professional Effectiveness.} Jossey-Bass.

\bibitem[Bandura, 1977]{bandura1977}
Bandura, A. (1977). Self-efficacy: Toward a unifying theory of behavioral change. \textit{Psychological Review, 84}(2), 191--215.

\bibitem[Bunge, 2015]{bunge2015}
Bunge, M. (2015). \textit{Causality and Mechanism in the Social Sciences.} Springer.

\bibitem[DellaVigna \& Linos, 2022]{dellavigna2022}
DellaVigna, S., \& Linos, E. (2022). RCTs to Scale: Comprehensive Evidence from Two Nudge Units. \textit{Econometrica, 90}(1), 81--116.

\bibitem[Dietz, 2026]{dietz2026}
Dietz, T. (2026). Building Actionable Theories: The Role of Causal Constructs. \textit{Behavioral Science Quarterly, 42}(3), 1--27.

\bibitem[Eisenhardt, 1999]{eisenhardt1999}
Eisenhardt, K. (1999). Strategy as Strategic Decision Making. \textit{Sloan Management Review.}

\bibitem[Epstein, 2023]{epstein2023}
Epstein, J. (2023). \textit{Generative Social Science 2.0.} MIT Press.

\bibitem[Fehr \& Fehr, 2024]{fehr2024}
Fehr, G., \& Fehr, A. (2024). \textit{Behavioral Change Modelling (BCM 2.2).} FehrAdvice Technical Report.

\bibitem[Giddens, 2000]{giddens2000}
Giddens, A. (2000). \textit{The Constitution of Society.} Polity Press.

\bibitem[Hummel \& Maedche, 2019]{hummel2019}
Hummel, D., \& Maedche, A. (2019). How effective is nudging? A quantitative review on the effect sizes and limits of empirical nudging studies. \textit{Journal of Behavioral and Experimental Economics, 80}, 47--58.

\bibitem[Pearl, 2009]{pearl2009}
Pearl, J. (2009). \textit{Causality: Models, Reasoning, and Inference.} Cambridge University Press.

\bibitem[Rousseau, 2018]{rousseau2018}
Rousseau, D. (2018). Evidence-Based Management Reconsidered. \textit{Academy of Management Perspectives.}

\bibitem[Simon, 1962]{simon1962}
Simon, H. (1962). The Architecture of Complexity. \textit{Proceedings of the American Philosophical Society, 106}(6), 467--482.

\bibitem[Weick, 1989]{weick1989}
Weick, K. (1989). Theory Construction as Disciplined Imagination. \textit{Academy of Management Review, 14}(4), 516--531.

\bibitem[Woodward, 2003]{woodward2003}
Woodward, J. (2003). \textit{Making Things Happen: A Theory of Causal Explanation.} Oxford University Press.

\end{thebibliography}

\end{document}
